
 
(TTĐN)-Người dân Can Lộc, Hà Tĩnh luôn tự hào về ``thần đồng'' toán học Phan Đình Diệu, một người con của quê hương. Năm 1965, khi mới 31 tuổi ông đã hoàn thành luận án tiến sĩ khoa học tại Đại học Tổng hợp quốc gia Moskva mang tên Lomonosov. Đối với những người nghiên cứu Toán học xứ Nghệ nói chung và Đại học Vinh nói riêng, GS là người anh cả.

Chiếc máy tính đầu tiên của Viện khoa học tính toán và điều khiển (Ảnh:TL)

Là người nghiên cứu toán học lâu năm, tôi có thể khẳng định GS Phan Đình Diệu là người có công lớn trong việc tạo ra ngành Tin học Việt Nam, từ việc đặt nền móng đến xây dựng định hướng phát triển. Là người có tâm với khoa học và rất yêu toán và tin học, anh cùng với những giáo sư đầu ngành như Tạ Quang Bửu, Hoàng  Tụy, Lê Văn Thiêm, Nguyễn Văn Đạo... bỏ không biết bao công sức để xây dựng ngành toán tin cho nước nhà.

Người cha mẫu mực

Nói về cuộc đời anh Diệu không thể không nói đến người đàn bà đứng sau. Đó là chị Văn Thị Xuân Hương, em gái của giáo sư Văn Như Cương. Hồi còn ở Nghệ An, anh Diệu và anh Cương cùng đi học với nhau. Năm 1957, sau khi tốt nghiệp đại học Sư phạm Hà Nội các anh đều về quê dạy học. Về quê, như lời giao ước, anh được anh Cương đưa ngay đến ``ra mắt'' cô em gái. Khi đó chị Xuân Hương mới 17 tuổi, kém anh Diệu 6 tuổi, thuộc loại xinh đẹp nhất trường, đang nằm viện vì cắt amidan. Nhiều chàng trong vùng dòm ngó, kể cả những bậc ``công tử'' quần là áo lượt, nhưng ``nàng'' chưa chấm ai.

%  Gia đình mừng sinh nhật tuổi 70 GS Phan Đình Diệu 
% (Ảnh: TLGĐ)

Nhưng anh Cương chỉ đưa ``ông em rể tương lai'' tới cổng bệnh viện rồi thả đó, kiểu ``mang con bỏ chợ'', rồi thủng thẳng: ``việc còn lại mi tự lo''. Lần đầu gặp ``thần đồng toán'', chị Hương, như có dòng điện chạy qua người khi thấy ông bạn của anh trai có gương mặt sáng sủa, trán cao, mắt đen...

Cụ thân sinh ra chị Hương thì đặc biệt mê anh Diệu, bạn con trai mình. Cụ bảo với gia đình, anh này rất được, cả nết lẫn chữ. Và ra một quyết định rất đồ Nghệ, ``không cho con gái lấy ai, ngoài anh Diệu!''. Thơ anh Diệu viết, cụ chép lại cho con gái, giờ vẫn bút tích còn trong sổ. Tán hộ, làm ``chân gỗ'' giúp con rể thì chưa từng thấy ai như cụ, một ông đồ Nghệ điển hình.

Năm 1961, anh Diệu đi nghiên cứu sinh ở Liên Xô, chị Hương về Vinh dạy học. Rồi chiến tranh mỗi người một nơi. Mãi mấy năm sau gia đình mới đoàn tụ khi anh Diệu về Hà Nội. Các con Quỳnh Dương, Hà Dương, rồi chú út Dương Hiệu ra đời.

Nói về chiều vợ, ít ai như anh Diệu. Chả là khi đó, Giáo sư đã có hai cô con gái là Quỳnh Dương và Hà Dương đều rất xinh xắn và hiền thục. Như nhiều gia đình xứ Nghệ khác, mọi người trong dòng họ muốn anh chị có con trai.
Anh Diệu nghiện thuốc lá nặng, phòng làm việc lúc nào cũng nghi ngút như cái lò gạch của Nam Cao. Ai đi công tác về cũng biếu anh thuốc lá, từ 555 đến Pall Mall, Malborro, Tam Đảo, Điện Biên, đủ loại tây ta.
Chị Hương khuyên chồng bỏ thuốc mãi không được nên ra giá: Nếu anh bỏ được thuốc lá, em sẽ đẻ con trai (?). Giáo sư hứa. Thật bất ngờ, cuối năm ấy, cháu trai Phan Dương Hiệu ra đời. Thế là anh đành phải thực hiện lời hứa bỏ thuốc! 


Mấy ngày đầu Giáo sư không chịu nổi, cứ cầm điếu thuốc lên, hút vài hơi rồi lại dụi tắt. Dần dần anh mang các bao thuốc bóc sẵn để khắp nơi trong nhà, trên giá sách, chỗ bàn ăn, văn phòng, nơi hay sờ tới nhất. Nhìn thấy thuốc nhưng không cầm. Sau ba tuần anh cai hẳn.

Đến giờ, các con anh chị đều đã thành đạt. Con gái là TS Quỳnh Dương, có 3 công chúa, hiện dạy ở Paris (Pháp). Cô thứ hai là Hà Dương, từng giành Huy chương Đồng Olympic Toán quốc tế năm 1990, tiến sỹ Toán hiện ở Hà Nội, ở cùng chung cư với bố mẹ tiện giúp đỡ các cụ. Cậu út Dương Hiệu, là giáo sư (Full Professor), hiện đang nghiên cứu và giảng dạy về mật mã bên Pháp, có vợ và con trai 7 tuổi thông minh còn hơn cả bố. Anh Diệu từng tâm sự, Dương Hiệu là nói ngược là Diệu- Hương, tên của anh Diệu và chị vợ là Hương ghép lại, kết quả của mối tình đẹp của anh chị.

% Vợ chồng Giáo sư Phan Đình Diệu năm 1964 (Ảnh: TLGĐ)

Người sáng lập Hội tin học Việt Nam

Năm 1975, trong một chuyến thực tập tại Pháp, ông đã được tiếp xúc với nhiều thành tựu hiện đại của ngành tin học trên thế giới.Từ đó, ông đã say mê tìm hiểu hai hướng phát triển mà ông cho là có triển vọng nhất và có thể ứng dụng và phát triển ở Việt Nam là vi tin học (trên cơ sở kỹ thuật vi xử lý và máy vi tính) và viễn tin học (trên cơ sở công nghệ viễn thông và mạng máy tính).

Hai năm sau, Viện Khoa học tính toán và điều khiển được chính thức thành lập, và ông được phân công làm Viện trưởng. Là người dự thảo kế hoạch và cũng là người quản lý, từ năm 1977 đến 1985, ông đã đưa Viện vượt qua nhiều khó khăn, trở ngại của buổi đầu hoạt động, xây dựng được một số hướng nghiên cứu chính về tin học.

Sau đó, ông làm Phó Viện trưởng Viện Khoa học Việt Nam (nay là Viện Khoa học và Công nghệ Việt Nam), Phó trưởng Ban thường trực Ban chỉ đạo Chương trình quốc gia về công nghệ thông tin (1993-1997), Thành viên Hội đồng Chính sách Khoa học và Công nghệ Quốc gia Việt Nam (từ năm 1992). Đi sang Pháp, thay vì dành tiền mua xe Peugeot về bán lấy lãi như nhiều người, anh mua rất nhiều linh kiện PC cho Viện. Nhớ có lần mưa lụt ngập hầm chứa máy tính ODRA, giáo sư đến thấy thảm cảnh đó mà rơi nước mắt vì biết bao tiền của mới mua được cái máy duy nhất thời đó ở miền Bắc.

Các seminar về automat hữu hạn, ngôn ngữ hình thức, lý thuyết mật mã, triết học và toán, do anh thuyết giảng hội trường luôn đông nghịt. Các bài giảng trừu tượng luôn được GS xứ Nghệ này đơn giản hóa cho dễ hiểu bởi anh cho rằng, cái đẹp Toán học nằm ở sự đơn giản và logic. Ông chính là người sáng lập và là Chủ tịch Hội tin học Việt Nam đầu tiên.

Khi đánh giá về GS Phan Đình Diệu, nhất là về các quan điểm xã hội sẽ có các ý kiến khác nhau, nhưng với chúng tôi, những người nghiên cứu toán biết tiếng Pháp gọi anh là Monsieur Dieu – anh Trời. Bài viết này như nén hương thơm kính cẩn nghiêng mình bên người Anh vừa tạ thế vào sáng qua, ngày 13/4 -``Ngày của Mẹ''- Anh đã trở về với Mẹ hiền, thọ 82 tuổi.

                                                       PGS.TS Nguyễn Quý Dy